\documentclass[11pt, a4paper]{article}

% ── PACKAGES ──────────────────────────────────────────────────
\usepackage[utf8]{inputenc}
\usepackage[T1]{fontenc}
\usepackage{geometry}
\geometry{margin=2.2cm}
\usepackage{hyperref}
\hypersetup{
    colorlinks=true,
    linkcolor=blue,
    urlcolor=blue,
    citecolor=blue,
    pdftitle={The GOT1/RUNX1 Transition Index: A Geometry-Derived
              2-Gene Attractor Depth Staging Tool for Clear Cell
              Renal Cell Carcinoma with Overall Survival Validation
              in TCGA-KIRC (n=532)},
    pdfauthor={Eric Robert Lawson / OrganismCore}
}
\usepackage{amsmath}
\usepackage{booktabs}
\usepackage{array}
\usepackage{longtable}
\usepackage{parskip}
\usepackage{microtype}
\usepackage{authblk}
\usepackage{titlesec}
\usepackage{fancyhdr}
\usepackage{xcolor}
\usepackage{float}
\usepackage{enumitem}
\usepackage{multirow}
\usepackage{mdframed}
\usepackage{tabularx}

% ── COLOURS ───────────────────────────────────────────────────
\definecolor{derivebox}{RGB}{235,245,255}
\definecolor{warningbox}{RGB}{255,245,200}
\definecolor{criticalbox}{RGB}{255,220,220}
\definecolor{findingbox}{RGB}{240,250,240}
\definecolor{infobox}{RGB}{235,245,255}
\definecolor{novelbox}{RGB}{248,240,255}
\definecolor{mechbox}{RGB}{245,255,245}
\definecolor{statsbox}{RGB}{250,245,255}
\definecolor{headerblue}{RGB}{30,80,160}

% ── HEADER / FOOTER ───────────────────────────────────────────
\pagestyle{fancy}
\fancyhf{}
\lhead{\footnotesize GOT1/RUNX1 Transition Index ---
       ccRCC Depth Staging Tool --- OrganismCore}
\rhead{\footnotesize 2026-03-07}
\rfoot{\small Page \thepage}
\renewcommand{\headrulewidth}{0.4pt}
\setlength{\headheight}{24pt}
\addtolength{\topmargin}{-10pt}

% ── SECTION FORMATTING ────────────────────────────────────────
\titleformat{\section}
  {\large\bfseries\color{headerblue}}
  {\thesection}{1em}{}[\titlerule]
\titleformat{\subsection}
  {\normalsize\bfseries}{\thesubsection}{1em}{}
\titleformat{\subsubsection}
  {\small\bfseries\itshape}{\thesubsubsection}{1em}{}

% ══════════════════════════════════════════════════════════════
\title{
    \vspace{-1em}
    \textbf{The GOT1/RUNX1 Transition Index:
    A Geometry-Derived 2-Gene Attractor Depth
    Staging Tool for Clear Cell Renal Cell
    Carcinoma with Overall Survival Validation
    in TCGA-KIRC ($n{=}532$)}\\[0.5em]
    \large Version 1.0\\[0.4em]
    \normalsize Waddington Landscape Derivation,
    Mechanistic Basis, Pre-Specified Prediction,
    Survival Confirmation, and IHC Translation
    Pathway for a Two-Gene Clinical Classifier
    of ccRCC Attractor Depth\\[0.3em]
    \small Supporting document for the OrganismCore
    RCC Series (Document RCC-S4-RA, 2026-03-07).
    Repository:
    \href{https://github.com/Eric-Robert-Lawson/attractor-oncology}
    {github.com/Eric-Robert-Lawson/attractor-oncology}
}

\author[1]{Eric Robert Lawson}
\affil[1]{OrganismCore \textbar{}
    \href{mailto:OrganismCore@proton.me}
    {OrganismCore@proton.me} \textbar{}
    ORCID: 0009-0002-0414-6544}

\date{March 7, 2026}

% ══════════════════════════════════════════════════════════════
\begin{document}
\maketitle
\thispagestyle{fancy}

% ── STATUS BOX ────────────────────────────────────────────────
\begin{mdframed}[backgroundcolor=warningbox,
                 linecolor=orange!70!black,
                 linewidth=1pt]
\textbf{STATUS:} Computational validation complete.
Pre-specified predictions confirmed in TCGA-KIRC
($n{=}532$, OS events $n{=}175$).
IHC H-score cut-points have \textbf{not yet been
established.} The Transition Index is currently
validated at the RNA expression level only.
IHC calibration (the remaining clinical step)
follows the same design as the FOXA1/EZH2 protocol
for breast cancer (OrganismCore CS-LIT-22,
\href{https://doi.org/10.5281/zenodo.18892788}
{DOI: 10.5281/zenodo.18892788}).
\end{mdframed}

\vspace{0.5em}
\tableofcontents
\newpage

% ══════════════════════════════════════════════════════════════
\section{Purpose and Scope}
% ══════════════════════════════════════════════════════════════

This document records the complete derivation,
mechanistic basis, pre-specified prediction,
computational survival validation, literature
check, and IHC translation pathway for the
GOT1/RUNX1 Transition Index (TI) as a 2-gene
attractor depth staging tool for clear cell
renal cell carcinoma (ccRCC).

The TI is defined as:

\begin{equation}
\text{TI} = \text{norm}(\text{GOT1})
           - \text{norm}(\text{RUNX1})
\label{eq:ti}
\end{equation}

\noindent where $\text{norm}(x)$ denotes the
z-score normalized RNA expression value of gene
$x$ across the tumour cohort.

\textbf{High TI} (GOT1-dominant):
proximal tubule identity retained;
shallow attractor position; better prognosis.

\textbf{Low TI} (RUNX1-dominant):
metabolic identity collapsed; chromatin false
attractor dominant; deep attractor position;
worse prognosis.

This document is intended for:

\begin{itemize}
    \item Researchers evaluating the ccRCC
    attractor depth framework and its
    clinical translation
    \item Pathologists and oncologists
    interested in a simple, low-cost
    depth stratification tool for ccRCC
    \item Institutions evaluating the TI
    for IHC calibration studies
    \item Groups pursuing biomarker
    stratification of belzutifan or other
    targeted therapy trials in ccRCC
\end{itemize}

% ══════════════════════════════════════════════════════════════
\section{How This Was Derived ---
Read This First}
\label{sec:derivation}
% ══════════════════════════════════════════════════════════════

\begin{mdframed}[backgroundcolor=derivebox,
                 linecolor=blue!50!black,
                 linewidth=1pt]
\textbf{NOT derived from survival data mining.
Derived from attractor geometry. The prediction
was locked before any survival analysis was run.}

\medskip
\textbf{Step 1 --- Geometry:}
Waddington landscape analysis of TCGA-KIRC
single-cell and bulk expression data identified
the primary attractor axis of ccRCC as a
transition from normal proximal tubule (PT)
metabolic identity to a chromatin-locked false
attractor state. The two anchor genes of this
axis --- GOT1 (falling with depth) and RUNX1
(rising with depth) --- were identified from
their depth correlation values:
$r(\text{GOT1}, \text{depth}) = -0.527$;
$r(\text{RUNX1}, \text{depth}) = +0.559$.
This was performed on the expression data
without reference to outcome data.

\medskip
\textbf{Step 2 --- Pre-specified prediction
(locked 2026-03-07 before Script 4 was run):}
\begin{itemize}[noitemsep]
    \item S4-P1: TI (as $-$TI continuous Cox)
    HR $> 1.5$, $p < 0.05$ in ccRCC OS
    \item S4-P2a: Q4 $<$ Q1 depth quartile
    OS (log-rank)
    \item S4-P2b: Q3+Q4 $<$ Q1+Q2 OS
    (log-rank)
    \item S4-P4: RUNX1, GOT1 confirmed in
    predicted OS directions
\end{itemize}
All four predictions were locked in the
pre-script document before the survival
dataset was merged.

\medskip
\textbf{Step 3 --- Confirmation:}
All four predictions confirmed. No additional
predictions were added after the run.
The statistics below are \textbf{confirmations
of a pre-specified prediction}, not the source
of it. This eliminates survival-endpoint fishing
and post-hoc framing as explanations for the
result. The before-document is publicly
timestamped in the OrganismCore repository.
\end{mdframed}

% ══════════════════════════════════════════════════════════════
\section{Biological Mechanism}
\label{sec:mechanism}
% ══════════════════════════════════════════════════════════════

\subsection{GOT1: The Proximal Tubule
Metabolic Identity Anchor}

GOT1 (Glutamic-Oxaloacetate Transaminase 1;
cytoplasmic aspartate aminotransferase, cAST)
is a central enzyme in amino acid metabolism
and the malate-aspartate shuttle. In normal
proximal tubule cells, GOT1 supports:

\begin{itemize}
    \item Aspartate synthesis and
    nucleotide biosynthesis
    \item Transfer of reducing equivalents
    (NADH) into mitochondria via the
    malate-aspartate shuttle
    \item TCA cycle anaplerosis through
    oxaloacetate supply
    \item Redox balance and mitochondrial
    respiratory function
\end{itemize}

In ccRCC, GOT1 falls continuously with
attractor depth ($r = -0.527$, $n = 534$
TCGA-KIRC tumours). This reflects the
metabolic collapse of proximal tubule
identity --- VHL loss drives HIF-mediated
glycolytic reprogramming, and the TCA cycle
enzymes that define normal PT cell identity
are progressively lost as the tumour
transitions deeper into the false attractor.

\textbf{Clinical laboratory note:}
GOT1 is the enzyme measured by the serum AST
(aspartate aminotransferase) test in standard
clinical panels. The serum De~Ritis ratio
(AST/ALT) is a known prognostic marker in RCC
at the \textit{systemic} level (elevated serum
ratio = worse OS, pooled HR 1.41, multiple
studies 2021--2024). The TI measures
\textit{tumour-intrinsic} GOT1 expression
at the RNA level --- a mechanistically distinct
quantity from the serum ratio, but the clinical
familiarity of GOT1/AST in the RCC context
is a translational asset.

\subsection{RUNX1: The False Attractor
Chromatin Lock Hub}

RUNX1 (Runt-Related Transcription Factor 1)
rises continuously with attractor depth
($r = +0.559$, $n = 534$). Its binding
partner CBFB co-rises ($r = +0.609$,
HR $= 2.60$ $[2.01$--$3.35]$, $p = 3.2
\times 10^{-13}$).

In the ccRCC false attractor, RUNX1 functions
as a transcriptional hub that:

\begin{itemize}
    \item Drives EZH2 and the broader
    Wall~2 chromatin lock (EZH2, DNMT3A,
    KDM1A, HDAC1 all rise with depth and
    co-vary with RUNX1)
    \item Maintains TGFBI-driven ECM
    adhesion (Wall~3), promoting the
    invasive programme in deep tumours
    \item Suppresses immune presentation
    through coupling to the IFI16$\to$B2M
    broken circuit in Q4 (Wall~4)
    \item Has been independently confirmed
    as an oncogenic driver in ccRCC by
    Cancer Research (2020): loss-of-function
    reduces tumour growth and improves
    survival in preclinical models
\end{itemize}

\textbf{Independent literature confirmation:}
Multiple published studies (PeerJ 2019/2020,
Cancer Research 2020, Immunity/Inflammation
and Disease 2024/2025, Scientific Reports 2025)
confirm RUNX1 as an independent adverse
prognostic marker in ccRCC (univariate Cox
HR $\approx 1.60$ $[1.31$--$1.97]$,
multivariate-confirmed). The present work
does not claim RUNX1's prognostic value
as novel. It claims the \textbf{mechanistic
framing of RUNX1 as the false attractor hub
and its ratio with GOT1 as a continuous
attractor depth index} is novel --- because
this framing and combination have no
prior publication.

\subsection{The Ratio as an Attractor
Position Meter}

The TI equation (Eq.~\ref{eq:ti}) captures
the balance between two opposing forces on
the proximal tubule identity axis:

\begin{center}
GOT1 (identity retained) vs
RUNX1 (identity locked off)
\end{center}

This is structurally identical to the
FOXA1/EZH2 ratio in breast cancer
(OrganismCore CS-LIT-1): FOXA1 opens
luminal identity; EZH2 silences it.
Their ratio measures which force is winning.
GOT1/RUNX1 measures the same contest on
the renal proximal tubule identity axis.

In both cases:
\begin{itemize}[noitemsep]
    \item The genes were identified from
    attractor geometry before any clinical
    data were examined
    \item The ratio was pre-specified
    before survival analysis
    \item The survival confirmation
    followed the pre-specification
    \item The IHC calibration is the
    remaining translational step
\end{itemize}

% ══════════════════════════════════════════════════════════════
\section{Pre-Specified Predictions and
Confirmation Status}
\label{sec:predictions}
% ══════════════════════════════════════════════════════════════

All predictions below were locked in the
OrganismCore RCC Script 4 pre-document
on 2026-03-07 before the survival dataset
was merged and Script 4 was executed.

\begin{mdframed}[backgroundcolor=findingbox,
                 linecolor=green!50!black,
                 linewidth=1.5pt]
\textbf{PREDICTION SCORECARD --- Script 4}

\medskip
\begin{tabular}{@{}p{5.5cm}p{2cm}p{5cm}@{}}
\toprule
\textbf{Prediction} &
\textbf{Verdict} &
\textbf{Result} \\
\midrule
S4-P1: $-$TI continuous Cox
HR $>$ 1.5, $p < 0.05$ &
\textbf{CONFIRMED} &
HR $= 6.94$ $[3.62$--$13.29]$,
$p = 5.09 \times 10^{-9}$,
$C = 0.627$ \\[0.5em]
S4-P2a: Q4 $<$ Q1 OS
(log-rank) &
\textbf{CONFIRMED} &
$p = 0.0001$;
Q1 median OS: not reached;
Q4 median OS: 1610 days \\[0.5em]
S4-P2b: Q3+Q4 $<$ Q1+Q2 OS
(log-rank) &
\textbf{CONFIRMED} &
$p = 9.60 \times 10^{-5}$;
shallow median: not reached;
deep median: 2227 days \\[0.5em]
S4-P4: RUNX1 high = worse OS &
\textbf{CONFIRMED} &
HR $= 1.58$, $C = 0.637$,
highest individual $C$-index \\[0.5em]
S4-P4: GOT1 low = worse OS &
\textbf{CONFIRMED} &
HR $= 0.67$, $p = 4.24
\times 10^{-9}$ \\[0.5em]
S4-P3: S5 depth $C >$
S2 depth $C$ &
\textbf{NOT CONFIRMED} &
$C_\text{S5} = 0.556$
vs $C_\text{S2} = 0.603$
(informative: composite
panel outperforms 2-gene TI
for survival prediction;
TI retains mechanistic and
clinical utility) \\
\bottomrule
\end{tabular}

\medskip
\textbf{Contradictions: 0.}
\end{mdframed}

% ══════════════════════════════════════════════════════════════
\section{Survival Analysis Results}
\label{sec:survival}
% ══════════════════════════════════════════════════════════════

\subsection{Dataset}

\begin{itemize}[noitemsep]
    \item \textbf{Dataset:} TCGA-KIRC
    (The Cancer Genome Atlas, Kidney Renal
    Clear Cell Carcinoma)
    \item \textbf{Expression samples:}
    $n = 534$ tumour samples
    \item \textbf{Survival merge:}
    $n = 532$ (after dropna on OS time
    and event)
    \item \textbf{OS events:} $n = 175$
    (32.9\% event rate)
    \item \textbf{OS range:} 2--4537 days
    \item \textbf{Endpoint:} Overall
    survival (OS) --- time to death from
    any cause
\end{itemize}

\subsection{Primary Result: Continuous Cox
Regression on TI}

The Transition Index was tested as a
continuous predictor of OS using Cox
proportional hazards regression.
The negative TI ($-$TI) was used so that
the hazard ratio is in the intuitively
correct direction (low TI = deep attractor
= higher hazard = HR $> 1$):

\begin{mdframed}[backgroundcolor=statsbox,
                 linecolor=purple!40!black,
                 linewidth=1pt]
\textbf{Primary survival result:}

\medskip
Cox regression on continuous $-$TI
(RUNX1-dominant direction):

\[
\text{HR} = 6.94 \quad
[95\%\,\text{CI:}\ 3.62\text{--}13.29]
\quad p = 5.09 \times 10^{-9}
\quad C\text{-index} = 0.627
\]

$n = 532$, events $= 175$.

\medskip
Kaplan-Meier median split on TI
(split at median TI $= -0.163$):
\begin{itemize}[noitemsep]
    \item TI-High (GOT1-dominant,
    $n = 266$): median OS
    \textbf{not reached}
    \item TI-Low (RUNX1-dominant,
    $n = 266$): median OS
    \textbf{1964 days}
    (${\approx}5.4$~years)
    \item Log-rank $p = 7.43
    \times 10^{-6}$
\end{itemize}
\end{mdframed}

\subsection{Quartile Depth Stratification}

Depth quartile assignment (Q1 = shallowest,
Q4 = deepest, $n = 133$ per quartile):

\begin{center}
\begin{tabular}{lccc}
\toprule
\textbf{Comparison} &
\textbf{Test} &
\textbf{$p$-value} &
\textbf{Verdict} \\
\midrule
Q4 vs Q1 (extreme quartiles) &
Log-rank &
$0.0001$ &
CONFIRMED \\[0.3em]
Q3+Q4 vs Q1+Q2 (combined) &
Log-rank &
$9.60 \times 10^{-5}$ &
CONFIRMED \\
\bottomrule
\end{tabular}
\end{center}

Q1 median OS: not reached.
Q4 median OS: 1610 days.
Q3+Q4 (deep) median OS: 2227 days.
Q1+Q2 (shallow) median OS: not reached.

\subsection{RUNX1 as the Dominant Individual
Transcription Factor Hub}

\begin{center}
\begin{tabular}{lcccc}
\toprule
\textbf{Gene} & \textbf{HR} &
\textbf{95\% CI} &
\textbf{$p$} &
\textbf{$C$} \\
\midrule
RUNX1 & 1.582 & 1.415--1.769 &
$8.08 \times 10^{-16}$ & 0.637 \\
CBFB  & 2.595 & 2.008--3.354 &
$3.20 \times 10^{-13}$ & 0.618 \\
IFI16 & 1.821 & 1.536--2.159 &
$5.04 \times 10^{-12}$ & 0.625 \\
EZH2  & 1.692 & 1.464--1.956 &
$1.08 \times 10^{-12}$ & 0.604 \\
KDM1A & 1.729 & 1.219--2.452 &
$0.0022$ & 0.521 \\
HDAC1 & 2.365 & 1.739--3.217 &
$4.19 \times 10^{-8}$ & 0.576 \\
GOT1  & 0.667 & 0.582--0.763 &
$4.24 \times 10^{-9}$ & 0.416 \\
OGDHL & 0.830 & 0.787--0.875 &
$6.92 \times 10^{-12}$ & 0.400 \\
SLC13A2 & 0.922 & 0.872--0.976 &
$0.0048$ & 0.460 \\
\bottomrule
\end{tabular}
\end{center}

RUNX1 achieves the highest individual
$C$-index ($0.637$) of all 21 genes tested.
In multivariate analysis including TI,
RUNX1 remains significant ($p = 0.002$),
confirming that RUNX1 carries independent
OS signal beyond the metabolic collapse
component captured by GOT1 alone.

\subsection{Multivariate Context}

\begin{center}
\begin{tabular}{lccl}
\toprule
\textbf{Model} & \textbf{Variable} &
\textbf{HR} & \textbf{$p$} \\
\midrule
\multirow{2}{*}{A: $-$TI + S2 depth}
& $-$TI & 4.84 & $0.0002$ \\
& S2 depth & 3.87 & $0.1886$ \\[0.3em]
\multirow{4}{*}{B: $-$TI + LOXL2 + TGFBI
+ RUNX1}
& $-$TI & 1.89 & $0.2432$ \\
& LOXL2  & 1.00 & $0.9413$ \\
& TGFBI  & 0.96 & $0.2607$ \\
& RUNX1  & 1.49 & $\mathbf{0.0020}$ \\
\bottomrule
\end{tabular}
\end{center}

\textit{Interpretation:} When TI and RUNX1
are entered together, RUNX1 carries residual
independent signal ($p = 0.002$). The TI
captures the metabolic identity axis
($-$GOT1 component); RUNX1 contributes
additionally as the transcription factor hub
component. For a simple clinical tool, TI
alone achieves $C = 0.627$ --- comparable
to established staging systems.

% ══════════════════════════════════════════════════════════════
\section{Literature Check --- Novelty
Confirmation}
\label{sec:litcheck}
% ══════════════════════════════════════════════════════════════

\subsection{What the literature has on
RUNX1 in ccRCC}

RUNX1 is a \textbf{confirmed prognostic
marker} in ccRCC. This is not a novel claim:

\begin{itemize}
    \item \textbf{PeerJ 2019/2020}
    (Qiu et al.): RUNX1 independently
    predicts OS in TCGA-KIRC,
    univariate HR $\approx 1.60$
    $[1.31$--$1.97]$, $p < 0.001$,
    multivariate confirmed
    \item \textbf{Cancer Research 2020}
    (AACR): RUNX1 is an oncogenic driver
    of ccRCC; loss-of-function reduces
    tumour growth in mouse models
    \item \textbf{Immunity, Inflammation
    and Disease 2024/2025}: RUNX1 as
    independent prognostic marker in ccRCC
    immune microenvironment; AUC up to
    $0.872$ for 1/3/5-year OS prediction
    \item \textbf{Scientific Reports 2025}:
    RUNX1 in angiogenesis-related
    multi-gene prognostic models in ccRCC
\end{itemize}

\textbf{Verdict on RUNX1 alone:}
Confirmed prognostic marker.
Not the novel contribution of this paper.

\subsection{What the literature has on
GOT1 in ccRCC}

\begin{itemize}
    \item \textbf{Frontiers in Oncology
    2024}: GOT1 in cancer metabolism ---
    critical role in amino acid metabolism,
    malate-aspartate shuttle, redox balance;
    GOT1 downregulated in ccRCC vs normal
    kidney in TCGA pan-cancer analysis
    \item \textbf{bioRxiv 2024}:
    Pan-cancer molecular signatures
    connecting aspartate metabolism (GOT1)
    to immune and metabolic features ---
    GOT1 downregulation in ccRCC is
    consistent with the metabolic identity
    collapse framework
    \item \textbf{De~Ritis ratio literature
    (serum AST/ALT, 2021--2024):}
    High preoperative serum De~Ritis ratio
    predicts worse OS in RCC (pooled
    HR 1.41 OS, 1.59 CSS; multiple papers).
    \textbf{Critical distinction:}
    The De~Ritis ratio measures
    \textit{serum enzyme levels} as a
    systemic hepatic/inflammatory stress
    marker. The TI measures
    \textit{tumour-intrinsic GOT1 RNA
    expression} as a cell identity anchor.
    These are mechanistically distinct
    constructs. Serum AST does not capture
    attractor depth.
\end{itemize}

\textbf{Verdict on GOT1 alone as
tumour-intrinsic attractor depth marker:}
Not in the literature as a standalone
ccRCC depth biomarker.

\subsection{What the literature does NOT have}

A systematic search for
\textit{``GOT1/RUNX1 ratio renal cell
carcinoma''} and \textit{``GOT1 RUNX1
combined prognostic index ccRCC''} confirms:

\begin{mdframed}[backgroundcolor=novelbox,
                 linecolor=purple!50!black,
                 linewidth=1pt]
\textbf{As of 2026-03-07, no published paper
reports a combined GOT1/RUNX1 gene ratio or
Transition Index as a prognostic tool in
renal cell carcinoma.}

The specific framing --- GOT1 as the metabolic
proximal tubule identity anchor and RUNX1 as
the false attractor chromatin lock hub, their
normalised difference as a continuous attractor
depth index with pre-specified OS survival
validation --- has no prior publication.

Multi-gene prognostic models exist in ccRCC
(lasso-Cox ferroptosis indices, immune gene
signatures, angiogenesis panels) but none
are mechanistically derived from Waddington
attractor geometry and none pair GOT1 and
RUNX1 as a 2-gene metabolic-vs-chromatin
axis.
\end{mdframed}

\textbf{Novel contribution of this paper:}
\begin{enumerate}
    \item The attractor geometry framing of
    GOT1 and RUNX1 as mechanistic opposites
    on the ccRCC identity axis
    \item The TI as a pre-specified,
    geometry-derived 2-gene continuous
    depth index
    \item The OS survival confirmation in
    $n = 532$ TCGA-KIRC with a
    pre-specified before-document
    (HR $= 6.94$, $p = 5.09 \times
    10^{-9}$, timestamps publicly verifiable)
    \item The IHC translation pathway
    (Section~\ref{sec:ihc})
    \item The treatment selection framework
    derived from TI quartile (Section~\ref{sec:treatment})
\end{enumerate}

% ══════════════════════════════════════════════════════════════
\section{The TI as a Treatment Selection
Framework}
\label{sec:treatment}
% ══════════════════════════════════════════════════════════════

The attractor depth quartile assigned by the
TI corresponds to distinct biology and distinct
drug logic. This is not a staging system in
the traditional TNM sense --- it is a
\textbf{mechanistic selection system} that
maps attractor position to therapeutic
vulnerability:

\begin{center}
\begin{tabularx}{\textwidth}{p{1.5cm}p{3.5cm}
                              p{3.5cm}p{3.5cm}}
\toprule
\textbf{TI / Quartile} &
\textbf{Dominant biology} &
\textbf{Appropriate drug logic} &
\textbf{What NOT to give} \\
\midrule
High TI Q1 (shallow) &
HIF/VEGF dominant; PT identity partially retained &
Belzutifan (HIF-2\textalphaantagonist); VEGF-targeted therapy;
cabozantinib + nivolumab &
--- \\[0.5em]
Mid TI Q2--Q3 &
HIF declining; chromatin lock emerging (EZH2,
DNMT3A rising) &
+ EZH2i (tazemetostat); + LOXL2i (emerging) &
--- \\[0.5em]
Low TI Q4 (deep) &
RUNX1/chromatin dominant; IFI16$\to$B2M
circuit broken; immune antigen presentation
absent; CD274 (PD-L1) flat/falling &
RUNX1i/CBFBi; + LOXL2i; + anti-B7-H3
(enoblituzumab); + $\alpha$KG supplementation;
Sequential: STING agonist
$\to$ HDACi $\to$ anti-PD-1 (CT-7 protocol) &
Anti-PD-L1 monotherapy (CD274 flat
in Q4; no target); Belzutifan alone
(HIF-2$\alpha$ not the dominant lock) \\
\bottomrule
\end{tabularx}
\end{center}

\subsection{The Belzutifan Resistance
Prediction (RCC-N1)}

The most immediately clinically urgent
application of the TI is patient stratification
for belzutifan (Welireg, FDA approved 2021
for advanced ccRCC in VHL mutation carriers).

\begin{mdframed}[backgroundcolor=criticalbox,
                 linecolor=red!60!black,
                 linewidth=1pt]
\textbf{NOVEL PREDICTION --- RCC-N1
(locked 2026-03-02, pre-script):}

RUNX1-high ccRCC patients (TI-Low / Q4) will
show inferior response to single-agent
belzutifan compared to RUNX1-low patients.

\medskip
\textbf{Mechanistic basis:}
Belzutifan targets HIF-2$\alpha$ (EPAS1).
In shallow tumours (Q1), EPAS1 is the
dominant attractor driver ($r(\text{EPAS1},
\text{depth}) < 0$; EPAS1 high = identity
retained = better OS --- confirmed in
Script 4: HR $= 0.67$, $p = 2.43
\times 10^{-18}$). As depth increases,
RUNX1 and the chromatin lock become
dominant. In Q4, the chromatin lock
(not HIF) is the primary attractor
mechanism. Belzutifan does not act on
RUNX1 or EZH2 or CBFB.

\medskip
\textbf{EPAS1 direction in Script 4:}
EPAS1 HR $= 0.666$ (high = better OS,
$p = 2.43 \times 10^{-18}$). This confirms
that EPAS1 elevation marks identity
retention (shallow attractor) in ccRCC.
The deeper the tumour, the lower EPAS1
--- the weaker the rational basis
for belzutifan.

\medskip
\textbf{Validation path:}
LITESPARK-005 or LITESPARK-013 trial data.
Stratify enrolled patients by RUNX1
expression (from archival tumour RNA).
Test response rate vs non-response vs
RUNX1 quartile.
This is testable from existing trial data
with a data access request to Merck or NCI.
No new patients required.
\end{mdframed}

% ══════════════════════════════════════════════════════════════
\section{IHC Translation Pathway}
\label{sec:ihc}
% ══════════════════════════════════════════════════════════════

\begin{mdframed}[backgroundcolor=warningbox,
                 linecolor=orange!70!black,
                 linewidth=1pt]
\textbf{STATUS:} The TI has been validated
at the RNA expression level. IHC H-score
cut-points have \textbf{not yet been
established.} The following section
describes the translation pathway --- not
the completed calibration.
\end{mdframed}

\subsection{The Principle}

The TI captures the relative balance between
GOT1 and RUNX1 protein expression in tumour
cells. If both proteins can be quantified
by IHC H-score in the same tissue section,
the ratio:

\[
\text{TI}_{\text{IHC}} =
\frac{\text{GOT1 H-score}}
     {\text{GOT1 H-score} + \text{RUNX1 H-score}}
\]

(or equivalent normalised difference)
would capture the same attractor balance
as the RNA-level TI, subject to calibration
against the RNA-validated cut-points.

\subsection{Antibody Availability}

\begin{center}
\begin{tabular}{p{2cm}p{3.5cm}p{3.5cm}p{3cm}}
\toprule
\textbf{Target} &
\textbf{Antibody / Clone} &
\textbf{Suppliers} &
\textbf{Notes} \\
\midrule
RUNX1 &
Clone EP206
(Leica Biosystems /
Novocastra) &
Leica, Abcam (EPR3099Y,
ab92336), CST \#4336,
GeneTex, Santa Cruz &
Nuclear staining.
Research-validated in
hematopoietic and solid
tumours. Requires
validation in FFPE
renal tissue. \\[0.5em]
GOT1 &
Multiple polyclonal
and monoclonal
antibodies available &
Abcam, CST, Sigma-Aldrich,
Proteintech,
Atlas Antibodies &
Cytoplasmic staining.
Used clinically in
neuroendocrine tumour
panels and hepatic
differentiation.
Requires validation
for ccRCC FFPE. \\
\bottomrule
\end{tabular}
\end{center}

\subsection{Calibration Study Design}

The calibration study required before
clinical deployment is directly analogous
to the FOXA1/EZH2 IHC calibration for
breast cancer (OrganismCore CS-LIT-22):

\begin{enumerate}
    \item \textbf{Cohort:} $n = 200$--$400$
    archived FFPE ccRCC cases with available
    RNA-seq or gene expression profiling data
    (TCGA tissue, institutional biobank, or
    TMA format)
    \item \textbf{Staining:} RUNX1 IHC +
    GOT1 IHC on serial sections or TMA cores
    \item \textbf{Scoring:} H-score
    (0--300) for each stain, cancer cell
    nuclei/cytoplasm only; scoring blinded
    to RNA-TI value and clinical outcome
    \item \textbf{Cut-point derivation:}
    Youden J optimisation of
    $\text{TI}_\text{IHC}$ against the
    RNA-level TI quartile assignment
    (Q1--Q4) as the ground truth reference
    \item \textbf{Validation:} Test
    IHC-TI cut-points against OS in
    the same cohort, with
    $\text{HR} > 1.5$ and $p < 0.05$
    as the confirmation threshold
\end{enumerate}

\textbf{Estimated resource requirement:}
Comparable to the FOXA1/EZH2 breast
cancer calibration: \$5,000--\$15,000
consumables, 3--6 months, depending on
cohort size and platform.

\textbf{Important caveat:} GOT1 IHC is
less routinely deployed in renal pathology
than FOXA1/EZH2 in breast pathology.
Additional antibody validation steps for
FFPE kidney tissue will be required before
the calibration study begins. This is a
tractable additional step, not a
fundamental obstacle.

% ══════════════════════════════════════════════════════════════
\section{Cross-Cancer Structural Significance:
The EZH2 Paradox in PRCC Type~2}
\label{sec:crosscancer}
% ══════════════════════════════════════════════════════════════

Script 4 produced an additional novel finding
in papillary RCC (PRCC) that provides
cross-cancer structural confirmation of the
attractor depth framework:

\begin{mdframed}[backgroundcolor=mechbox,
                 linecolor=green!40!black,
                 linewidth=1pt]
\textbf{THE EZH2 PARADOX IN PRCC TYPE~2
(novel, TCGA-KIRP, confirmed 2026-03-07):}

\medskip
In PRCC Type 2 ($n = 142$ with GDC subtype
confirmation, $n_\text{events} = 31$):

\begin{itemize}[noitemsep]
    \item EZH2 univariate Cox:
    HR $= 1.85$, $p = 0.0066$
    (high EZH2 $=$ worse OS)
    \item EZH2 multivariate (adjusting
    for MKI67): HR $= 0.19$,
    $p = 0.0041$
    (\textbf{high EZH2 $=$ better OS}
    within MKI67-high population)
\end{itemize}

\medskip
\textbf{Interpretation:} EZH2-low +
MKI67-high is the most dangerous
phenotype in PRCC Type~2. EZH2-high
tumours within the MKI67-high population
are epigenetically committed --- less
plastic, less invasive --- and therefore
paradoxically associated with better OS
within this proliferative subgroup.

\medskip
\textbf{This is the same paradox mechanism
confirmed in breast cancer TNBC
(OrganismCore CS-LIT-10, GSE25066,
$n = 508$):}
Attractor depth (EZH2-high = more
committed) predicts short-term
chemotherapy sensitivity AND long-term
progression risk, with the two arms
separable by proliferative context
(MKI67 in PRCC; chemotherapy response
in TNBC).

\medskip
\textbf{Same mechanism. Different cancer.
Different dataset. Different platform.}
This is cross-cancer structural
confirmation of the attractor depth
framework.
\end{mdframed}

% ══════════════════════════════════════════════════════════════
\section{Honest Limits}
\label{sec:limits}
% ══════════════════════════════════════════════════════════════

\begin{enumerate}

    \item \textbf{Single dataset.}
    The survival validation is in
    TCGA-KIRC only ($n = 532$).
    External replication in an
    independent ccRCC cohort with
    OS data and RNA-seq is required
    before clinical claims can
    be made. The TCGA result is
    strong (HR $= 6.94$, pre-specified)
    but single-dataset.

    \item \textbf{RNA-level only.}
    No IHC H-score cut-points exist.
    The clinical tool is not deployable
    from the current data. The
    calibration study is required.

    \item \textbf{All drug predictions
    are preclinical hypotheses.}
    The treatment selection framework
    in Section~\ref{sec:treatment}
    is geometry-derived and literature-
    supported but has not been
    prospectively tested. No patients
    have been treated according to
    this framework.

    \item \textbf{GOT1 IHC validation
    in renal tissue.}
    Unlike FOXA1/EZH2 in breast cancer
    (both antibodies are in routine
    clinical use in breast pathology),
    GOT1 IHC requires validation in FFPE
    renal tissue before the calibration
    study can begin. RUNX1 IHC is
    commercially available with multiple
    validated clones but is not routine
    in renal pathology. This is an
    additional antibody validation step
    not present in the breast cancer
    protocol.

    \item \textbf{The PRCC EZH2 paradox
    finding is low-power.}
    $n_\text{events} = 31$ in PRCC
    Type~2 is sufficient to detect the
    multivariate reversal but requires
    replication in a larger PRCC cohort.

\end{enumerate}

% ══════════════════════════════════════════════════════════════
\section{Repository and Timestamp
Verification}
\label{sec:repo}
% ══════════════════════════════════════════════════════════════

All pre-specified predictions, analysis
scripts, and reasoning artifacts are publicly
available in the OrganismCore GitHub repository
with commit-level timestamps that predate the
survival analysis run:

\begin{center}
\href{https://github.com/Eric-Robert-Lawson/attractor-oncology}
{\texttt{github.com/Eric-Robert-Lawson/attractor-oncology}}
\end{center}

Relevant documents:
\begin{itemize}[noitemsep]
    \item \texttt{Cancer\_Research/RCC/CCRCC/CCRCC\_false\_attractor\_v5.md}
    --- Script 5 pre-locked novel predictions
    (N1--N7), including the TI definition
    and the belzutifan resistance prediction
    \item \texttt{Cancer\_Research/RCC/Analysis\_After\_BRCA/RCC\_N5\_Validated\_Cat\_A.md}
    --- Script 4 survival analysis reasoning
    artifact (post-run)
    \item \texttt{Cancer\_Research/RCC/Analysis\_After\_BRCA/RCC\_Master\_Artifact\_From\_Lit\_Checks.md}
    --- Complete RCC series master findings
    \item \texttt{BRCA\_30\_30\_coherence.md}
    --- Cross-series coherence argument
\end{itemize}

% ══════════════════════════════════════════════════════════════
\section{Summary}
\label{sec:summary}
% ══════════════════════════════════════════════════════════════

\begin{mdframed}[backgroundcolor=findingbox,
                 linecolor=green!50!black,
                 linewidth=1.5pt]

The GOT1/RUNX1 Transition Index is a
geometry-derived 2-gene attractor depth
staging tool for clear cell renal cell
carcinoma.

\medskip
\textbf{What it is:}
$\text{TI} = \text{norm}(\text{GOT1}) -
\text{norm}(\text{RUNX1})$.
GOT1 = proximal tubule metabolic identity
anchor (falling with depth).
RUNX1 = false attractor chromatin lock
hub (rising with depth).
Their normalised difference captures
attractor position continuously.

\medskip
\textbf{What the data show:}
In TCGA-KIRC ($n = 532$, pre-specified):
Cox HR $= 6.94$ $[3.62$--$13.29]$,
$p = 5.09 \times 10^{-9}$, $C = 0.627$.
TI-Low median OS: 1964 days vs
TI-High median OS: not reached
($p = 7.43 \times 10^{-6}$).
Quartile stratification: Q4 vs Q1
$p = 0.0001$. 20/21 individual gene OS
directions confirmed.

\medskip
\textbf{What is novel:}
The GOT1/RUNX1 ratio as an attractor-
geometry-derived depth index is not in
the literature. The pre-specified,
geometry-first derivation is documented
with public timestamps. The individual
prognostic value of RUNX1 in ccRCC is
confirmed but not the novel contribution.

\medskip
\textbf{What remains:}
IHC calibration study in archival FFPE
ccRCC tissue. The clinical tool is not
deployable from current data.
The calibration study is the remaining step.

\medskip
\textbf{The single sentence:}
A two-gene metabolic-vs-chromatin ratio,
derived from attractor geometry before
any survival data were examined, stratifies
clear cell renal cell carcinoma patients
by depth of attractor commitment and
predicts overall survival with HR $= 6.94$
in a pre-specified, independently
confirmable analysis.

\end{mdframed}

% ── AUTHOR / AFFILIATION ──────────────────────────────────────
\vspace{1em}
\noindent\rule{\textwidth}{0.4pt}

\noindent\textbf{Author:} Eric Robert Lawson\\
\textbf{Affiliation:} OrganismCore
(Independent Research)\\
\textbf{Email:}
\href{mailto:OrganismCore@proton.me}
{OrganismCore@proton.me}\\
\textbf{ORCID:}
\href{https://orcid.org/0009-0002-0414-6544}
{0009-0002-0414-6544}\\
\textbf{Repository:}
\href{https://github.com/Eric-Robert-Lawson/attractor-oncology}
{github.com/Eric-Robert-Lawson/attractor-oncology}\\
\textbf{Date:} March 7, 2026\\
\textbf{Version:} 1.0 (Initial release)

\end{document}
